\chapter{Flat Feet and Fallen Arches}

\section*{Definition and Overview}
Flat feet and fallen arches describe the same underlying finding: the arch of the foot has lowered or collapsed, so the sole makes more contact with the ground than usual. "Flat feet" (pes planus) usually describes the shape of the foot, while "fallen arches" often refers to the arch lowering over time, particularly in adulthood. Most flat feet are flexible and painless and need no treatment, but a fallen arch that appears suddenly, is painful, or affects only one foot can signal a problem such as a torn tendon and warrants assessment. This chapter explains the different types of flat feet, their causes, and when they matter.

\section*{Epidemiology}
Flat feet are common at all ages. Most infants and toddlers have flat feet, and arches develop naturally by school age in the majority of children. In adults, around 20--30\% have flat feet to some degree, and most are asymptomatic. Adult-acquired flat foot from posterior tibial tendon dysfunction is much less common but important, affecting mainly middle-aged and older adults, women more than men, and people who are overweight or have diabetes or hypertension. Painful flat feet affect a minority overall, but they are a frequent reason for podiatry and orthopaedic consultation.

\section*{Aetiology and Pathophysiology}
\begin{itemize}
    \item \textbf{Physiological flat feet in children:} arches develop as the foot's muscles and connective tissue strengthen with weight bearing; most children outgrow flat feet
    \item \textbf{Flexible flat feet:} the arch flattens on standing but returns when weight is lifted; usually inherited, associated with flexible joints, and typically harmless
    \item \textbf{Adult-acquired flat foot:} posterior tibial tendon dysfunction, in which the tendon supporting the arch weakens, stretches, or tears, causing a painful, progressive collapse
    \item \textbf{Rigid flat feet:} the arch is flat even without weight bearing, often from tarsal coalition (fusion of foot bones) or arthritis
    \item \textbf{Contributing factors:} obesity, prolonged standing, pregnancy, diabetes, high-impact sport, and previous foot injury
    \item \textbf{Effects of collapse:} the foot rolls inward (overpronation), altering forces through the ankle, knee, hip, and spine
    \item \textbf{When it matters:} pain, stiffness, swelling, or a change in the shape of one foot are signs the flat foot is not just a shape but a problem
\end{itemize}

\section*{Clinical Features}
\begin{itemize}
    \item The inner arch touches or nearly touches the ground when standing
    \item The foot and heel tilt inward when viewed from behind (overpronation)
    \item No pain or disability in the majority of people
    \item Pain in the arch, heel, inner ankle, or calf, worse with standing, walking, or running
    \item Swelling along the inner ankle and difficulty standing on tiptoes in tendon dysfunction
    \item A rigid, flat foot with limited movement, sometimes painful in adolescents
    \item Rapid or uneven wear on the inner soles of shoes
    \item Knock-knees, shin pain, or knee and back discomfort from altered gait
    \item Sudden or one-sided collapse of an arch in an adult is a warning sign
\end{itemize}

\section*{Differential Diagnosis}
Foot and ankle pain needs sorting into its possible causes.
\begin{itemize}
    \item Posterior tibial tendon dysfunction, the main cause of adult-acquired flat foot
    \item Plantar fasciitis, which causes heel pain and often coexists with flat feet
    \item Tarsal coalition in children and adolescents with rigid, painful flat feet
    \item Arthritis, including osteoarthritis and inflammatory arthritis such as rheumatoid disease
    \item Stress fractures of the foot or ankle
    \item Achilles tendinopathy and other tendon problems
    \item Peripheral neuropathy in diabetes, which can weaken foot muscles
\end{itemize}

\section*{When to Seek Medical Care}
Painless flat feet need no treatment. See a doctor or podiatrist for foot pain, swelling, stiffness, a sudden change in arch shape (especially one-sided), difficulty walking, or flat feet that limit activity.

\textbf{Call 000 (or local emergency services) for:}
\begin{itemize}
    \item Sudden severe foot or ankle pain with swelling after injury
    \item An inability to bear weight
    \item Sudden calf swelling, pain, or redness, which may indicate a blood clot
    \item A pale, cold, or numb foot
\end{itemize}

\section*{Diagnosis and Investigations}
\begin{itemize}
    \item \textbf{Examination:} observation of the arch standing and sitting, the heel position from behind, and the walking pattern
    \item \textbf{Tiptoe test:} rising onto the toes should restore the arch; failure suggests posterior tibial tendon dysfunction
    \item \textbf{Footwear inspection:} wear patterns reveal how the foot strikes the ground
    \item \textbf{X-rays:} standing views of the feet show alignment and exclude arthritis and coalition
    \item \textbf{Ultrasound or MRI:} for suspected tendon tears, ligament injury, or stress fractures
    \item \textbf{Blood tests:} for inflammatory arthritis when suspected
\end{itemize}

\section*{Management}
\begin{itemize}
    \item \textbf{Reassurance:} flexible, painless flat feet, especially in children, require only supportive shoes
    \item \textbf{Footwear:} stable shoes with arch support and a firm heel counter are the foundation of management
    \item \textbf{Orthotics:} arch supports and insoles relieve symptoms in many adults
    \item \textbf{Physiotherapy:} strengthening of the posterior tibial tendon, calf muscles, and foot intrinsic muscles, with stretches
    \item \textbf{Pain relief:} ice and anti-inflammatory medicines for flare-ups
    \item \textbf{Weight management:} reducing excess weight eases the load on the arches
    \item \textbf{Tendon dysfunction:} early treatment with rest, orthotics, physiotherapy, and sometimes a walking boot; surgery if the tendon tears or the foot stiffens
    \item \textbf{Surgery:} reserved for rigid painful flat feet or failed conservative treatment
\end{itemize}

\section*{Complications}
\begin{itemize}
    \item Chronic foot, ankle, knee, hip, and back pain from malalignment
    \item Plantar fasciitis and Achilles problems
    \item Progressive deformity and stiffness from untreated tendon dysfunction
    \item Reduced mobility, activity, and quality of life
    \item Increased falls risk in older adults
\end{itemize}

\section*{Prognosis and Outcomes}
Most flat feet cause no trouble and require no treatment. Symptomatic flat feet usually improve substantially with supportive footwear, orthotics, and physiotherapy, and surgery is rarely needed. Adult-acquired flat foot from tendon dysfunction has a much better outcome with early treatment, which is why a newly fallen, painful arch should be assessed promptly. With appropriate care, most people with flat feet remain active and pain-free.

\section*{Prevention and Self-Care}
\begin{itemize}
    \item Wear supportive shoes that fit well, including around the house and at work
    \item Use arch supports if they help
    \item Strengthen the feet and calves with regular exercises
    \item Stretch the calf muscles and Achilles tendon
    \item Maintain a healthy weight
    \item Rest and ice the feet after prolonged standing or activity
    \item Replace worn shoes regularly
    \item See a podiatrist or doctor early for pain or a change in foot shape
\end{itemize}

\section*{Sources and Further Reading}
The following reputable sources provide further information for healthcare students and patients seeking reliable, current advice about flat feet and fallen arches.
\begin{itemize}
    \item Healthdirect Australia. \textit{Flat feet.} https://www.healthdirect.gov.au/flat-feet
    \item Australian Podiatry Association. \textit{Flat feet.} https://www.podiatry.org.au/
    \item Better Health Channel. \textit{Flat feet.} https://www.betterhealth.vic.gov.au/
    \item NHS. \textit{Flat feet.} https://www.nhs.uk/conditions/flat-feet/
    \item Mayo Clinic. \textit{Flatfeet.} https://www.mayoclinic.org/
    \item MedlinePlus. \textit{Flat feet.} https://medlineplus.gov/
\end{itemize}
